\documentclass[12pt,a4paper]{article}

\usepackage[ngerman]{babel}
\usepackage[utf8]{inputenc}
\usepackage[T1]{fontenc}
\usepackage{geometry}
\geometry{a4paper,margin=2.5cm}
\usepackage{hyperref}
\usepackage{graphicx}
\usepackage{listings}
\usepackage{xcolor}
\usepackage{booktabs}
\usepackage{tabularx}
\usepackage{amsmath,amssymb}
\usepackage{enumitem}
\usepackage{fancyhdr}
\usepackage{caption}
\usepackage{tikz}
\usepackage{lmodern}
\renewcommand{\familydefault}{\sfdefault}

\usetikzlibrary{positioning,fit,backgrounds,arrows.meta,shapes.geometric,chains}

\linespread{1.15}

\definecolor{codegreen}{rgb}{0,0.6,0}
\definecolor{codegray}{rgb}{0.5,0.5,0.5}
\definecolor{codepurple}{rgb}{0.58,0,0.82}
\definecolor{bg}{rgb}{0.96,0.96,0.96}

\lstdefinelanguage{TypeScript}{
  keywords={abstract, any, as, async, await, boolean, break, case, catch, class, console, const, continue, debugger, declare, default, delete, do, else, enum, export, extends, false, finally, for, function, get, if, implements, import, in, instanceof, interface, is, keyof, let, module, namespace, never, new, null, number, object, package, private, protected, public, readonly, require, return, set, static, string, super, switch, symbol, this, throw, true, try, type, typeof, undefined, var, void, while, with, yield},
  keywordstyle=\color{blue}\bfseries,
  ndkeywords={class, export, boolean, throw, implements, import, this},
  ndkeywordstyle=\color{darkgray}\bfseries,
  identifierstyle=\color{black},
  sensitive=true,
  comment=[l]{//},
  morecomment=[s]{/*}{*/},
  commentstyle=\color{codegreen}\ttfamily,
  stringstyle=\color{codepurple}\ttfamily,
  morestring=[b]',
  morestring=[b]"
}

\lstdefinelanguage{yaml}{
  keywords={true,false,null,y,n,version,services,build,ports,volumes,environment,depends_on},
  keywordstyle=\color{blue}\bfseries,
  basicstyle=\ttfamily\footnotesize,
  sensitive=false,
  comment=[l]{\#},
  commentstyle=\color{codegreen}\ttfamily,
  stringstyle=\color{codepurple}\ttfamily,
  morestring=[b]',
  morestring=[b]"
}

\lstset{
  backgroundcolor=\color{bg},   
  commentstyle=\color{codegreen},
  keywordstyle=\color{blue},
  numberstyle=\tiny\color{codegray},
  stringstyle=\color{codepurple},
  basicstyle=\ttfamily\footnotesize,
  breakatwhitespace=false,
  breaklines=true,
  captionpos=b,
  keepspaces=true,
  numbers=left,
  numbersep=5pt,
  showspaces=false,
  showstringspaces=false,
  showtabs=false,
  tabsize=2,
  frame=single
}

\hypersetup{
  colorlinks=true,
  linkcolor=blue,
  filecolor=magenta,
  urlcolor=cyan,
  pdftitle={Flappy Bird - Reimagined},
  pdfauthor={Alexander Engelhardt}
}

\setlength{\headheight}{15pt}
\pagestyle{fancy}
\fancyhf{}
\rhead{Flappy Bird -- Reimagined}
\lhead{Web-Applikationen SS26}
\rfoot{\thepage}

\begin{document}

% Deckblatt
\begin{titlepage}
  \begin{center}
    \vspace*{0.5cm}
    {\LARGE\bfseries Hochschule Konstanz (HTWG)}\\[0.3cm]
    {\large Fakultät Informatik -- Web-Applikationen SS26}\\[1.5cm]

    \rule{\linewidth}{0.5mm}\\[0.4cm]
    {\LARGE\bfseries Flappy Bird -- Reimagined}\\[0.2cm]
    \rule{\linewidth}{0.5mm}\\[1.2cm]

    {\Large\itshape Abschlussprojekt \& Ausarbeitung (Meilenstein 4)}\\[2.5cm]
  \end{center}

  \vfill

  \noindent
  \begin{minipage}[t]{0.48\textwidth}
    {\large\bfseries Ein Projekt von:}\\
    Alexander Engelhardt\\[1em]
    {\large\bfseries Matrikelnummer:}\\
    313291\\[1em]
    {\large\bfseries Studiengang:}\\
    Angewandte Informatik\\[1em]
    {\large\bfseries Fachsemester:}\\
    5
  \end{minipage}
  \hfill
  \begin{minipage}[t]{0.48\textwidth}\raggedleft
    {\large\bfseries Prüfer:}\\
    Pascal Laube\\[2em]
    {\large\bfseries Abgabedatum:}\\
    07. August 2026
  \end{minipage}

  \vspace*{0.5cm}
  \end{titlepage}

% Verzeichnisse
\newpage
\pagenumbering{roman}
\tableofcontents
\newpage
\listoftables
\listoffigures
\lstlistoflistings
\newpage
\pagenumbering{arabic}

% Kap. 1: Einleitung
\section{Einleitung und Problemstellung}

\subsection{Historischer Kontext des Ur-Spiels}
Das im Jahr 2013 vom vietnamesischen Entwickler Dong Nguyen veröffentlichte Spiel \textit{Flappy Bird} stellte ein phänomenales Zeitdokument der mobilen Videospielgeschichte dar. Aufgrund seiner rudimentären Steuerung -- ein einzelner Antipp-Impuls zur Überwindung der Gravitation -- und seines extrem hohen Schwierigkeitsgrads erreichte das Spiel innerhalb kürzester Zeit hunderte Millionen Downloads. 

Trotz dieses kommerziellen Phänomens wies das ursprüngliche Spiel gravierende technische und architektonische Mängel auf, die aus heutiger Sicht moderner Softwareentwicklung nicht mehr tragbar sind:
\begin{enumerate}
  \item \textbf{Fehlende Serveranbindung:} Sämtliche Spielstände wurden ausschließlich lokal auf dem Endgerät gespeichert, was die Verifikation von Rekorden unmöglich machte.
  \item \textbf{Anfälligkeit für Cheating:} Da keine serverseitige Validierung stattfand, konnten lokale Spielstände mit einfachen Speicher-Editoren manipuliert werden.
  \item \textbf{Fehlendes Benutzer-Management:} Es existierte keine Möglichkeit, Benutzerkonten anzulegen, individuelle Statistiken einzusehen oder sich authentifiziert in globalen Ranglisten zu messen.
  \item \textbf{Statisches visuelles Design:} Das Spiel bot keinerlei Anpassungsmöglichkeiten für Barrierefreiheit oder unterschiedliche Lichtverhältnisse (z.\,B. Darkmode/Lightmode).
\end{enumerate}

\subsection{Projektziel: Flappy Bird -- Reimagined}
Im Rahmen der Lehrveranstaltung \textit{Web-Applikationen} im Sommersemester 2026 an der Hochschule Konstanz (HTWG) wurde das Projekt **Flappy Bird -- Reimagined** ins Leben gerufen. Ziel war es, die ikonische Kernmechanik des Spiels aufzugreifen und in eine moderne, modulare, typsichere Full-Stack-Webanwendung zu überführen.

Das Projekt verfolgt hierbei folgende Hauptzielsetzungen:
\begin{itemize}
  \item \textbf{Plattformunabhängigkeit:} Vollständige Lauffähigkeit im Browser ohne zusätzliche Plugins oder native Installationen.
  \item \textbf{Typsicherheit über den gesamten Stack:} Durchgängiger Einsatz von TypeScript im Frontend und Backend zur Minimierung von Laufzeitfehlern.
  \item \textbf{Moderne Authentifizierung:} Implementierung einer zustandslosen Session-Verwaltung mittels JSON Web Tokens (JWT) und sicherer Passwort-Verschlüsselung via Bcrypt.
  \item \textbf{Persistente Highscore-Verwaltung:} Aufbau einer relationalen SQLite-Datenbank zur verlässlichen Speicherung von Bestenlisten.
  \item \textbf{Containerisiertes Deployment:} Vollständige Entkopplung der Systemkomponenten mittels Docker und Docker-Compose für ein reproduzierbares Deployment auf beliebigen Zielsystemen.
\end{itemize}

\subsection{Narrativ und Storytelling}
Zur Steigerung der User-Retention und Verbesserung der UI/UX-Experience wurde das Spiel um ein eigenständiges Lore-Konzept erweitert. Auf der Landingpage (\texttt{Index.tsx}) wird der Spieler in die Welt von Flappy eingeführt:

\begin{quote}
„Long ago, a small yellow bird lived in a quiet world above endless green pipes. Unlike the other birds, Flappy was never meant to stay in a cage or rest on the ground. It dreamed of the open sky and of proving that even the smallest wings could overcome impossible obstacles. One day, the bird escaped captivity and began its journey through the pipe-filled land, where every flight became a test of courage, timing, and endurance.“
\end{quote}

\subsection{Zielgruppenanalyse und Anwendungsfälle}
Die Zielgruppe der Anwendung teilt sich in zwei primäre Segmente auf:

\textbf{1. Competitive Casual-Gamer:}
Nutzer, die reaktionsbasierte, kurze Spielrunden zur Zerstreuung suchen, sich jedoch durch die globale Top-10-Bestenliste zu kontinuierlichen Verbesserungsschritten anspornen lassen. Für diese Zielgruppe stehen Performanz, geringe Latenz im Game-Loop und sofortiges Highscore-Feedback im Vordergrund.

\textbf{2. Gelegenheitsnutzer \& Familien:}
Anwender, die eine barrierefreie Webanwendung ohne Barrieren suchen. Durch den integrierten Light/Dark-Theme-Toggle wird das Spiel an unterschiedliche Umgebungsbeleuchtungen angepasst.

% Kap. 2: Technologischer Deep Dive
\section{Technologie-Stack und Werkzeugkette}

Der gewählte Stack basiert auf bewährten Industriestandards des JavaScript/TypeScript-Ökosystems. Tabelle~\ref{tab:tech-stack-detail} bietet eine komponentenweise Aufschlüsselung der eingesetzten Bibliotheken.

\begin{table}[htbp]
  \centering
  \begin{tabularx}{\textwidth}{lXl}
    \toprule
    \textbf{Komponente} & \textbf{Technologie / Bibliothek} & \textbf{Version} \\ \midrule
    Client-Framework & React (SPA) & 19.x \\
    Language Standard & TypeScript & 5.x \\
    Build Tool / Bundler & Vite & 6.x \\
    Client-Routing & React Router DOM & 7.x \\
    HTTP-Client & Axios & 1.7.x \\
    Backend-Server & Node.js mit Express.js & 4.21.x \\
    Passwort-Verschlüsselung & Bcrypt & 5.1.x \\
    Session-Sicherheit & JsonWebToken (JWT) & 9.0.x \\
    Relationales DBMS & SQLite (\texttt{better-sqlite3}) & 11.x \\
    Containerisierung & Docker \& Docker-Compose & 2.x \\ \bottomrule
  \end{tabularx}
  \caption{Detaillierte Analyse des Technologie-Stacks}
  \label{tab:tech-stack-detail}
\end{table}

\subsection{Frontend: React 19, Vite und TypeScript}
Die Frontend-Architektur setzt auf \textbf{React 19} in Kombination mit \textbf{Vite}. Die Entscheidung für Vite anstelle klassischer Tools wie Webpack begründet sich durch das extrem schnelle Hot Module Replacement (HMR) auf Basis von nativen ES-Modulen. TypeScript fungiert als verbindendes Element über das gesamte Frontend. Durch strikte Typdefinitionen (z.\,B. für Pipe-Objekte oder Formular-Eigenschaften in \texttt{types.ts}) werden Refactorings und Zustandstransformationen bereits zur Kompilierzeit abgesichert.

\subsection{Backend: Node.js, Express und SQLite}
Als Server-Lösung kommt \textbf{Node.js} mit dem \textbf{Express.js}-Framework zum Einsatz. Die Asynchronität von Node.js prädestiniert das Framework für I/O-intensive Webapplikationen. Als Datenbank wird \textbf{SQLite} über das Node-Modul \texttt{better-sqlite3} eingebunden. Im Gegensatz zu herkömmlichen asynchronen SQLite-Treibern arbeitet \texttt{better-sqlite3} synchron, was den Codenavigationsaufwand reduziert und bei lokalen Lese-/Schreibzugriffen eine extrem hohe Durchsatzrate garantiert.

\subsection{Sicherheit: Bcrypt und JWT}
Zur Passwortsicherung wird der Bcrypt-Algorithmus genutzt. Durch die Inkorporierung eines variablen Salt-Faktors (standardmäßig 10 Runden) wird das System gegen Rainbow-Table-Angriffe geschützt. Das Session-Management erfolgt über JSON Web Tokens (JWT). Nach erfolgreicher Authentifizierung stellt der Server ein kryptographisch signiertes Token aus, welches vom Client bei geschützten Anfragen im \texttt{Authorization: Bearer <token>}-Header mitgeführt werden muss.

% Kap. 3: Systemarchitektur & Entwurf
\section{Systemarchitektur und Systementwurf}

\subsection{Drei-Schichten-Architektur}
Die Webapplikation basiert auf einer entkoppelten Drei-Schichten-Architektur. Abbildung~\ref{fig:arch-diag} visualiert das Zusammenspiel zwischen Client-Browser, Express-Backend und SQLite-Datenbank.

\begin{figure}[htbp]
  \centering
  \begin{tikzpicture}[
    font=\sffamily\small,
    node distance=1.5cm and 1.8cm,
    box/.style={draw=blue!70!black, fill=blue!5, rounded corners=4pt, align=center, inner sep=8pt, minimum width=3.8cm, minimum height=1.2cm},
    serverbox/.style={draw=green!60!black, fill=green!5, rounded corners=4pt, align=center, inner sep=8pt, minimum width=3.8cm, minimum height=1.2cm},
    dbbox/.style={draw=orange!80!black, fill=orange!5, rounded corners=4pt, align=center, inner sep=8pt, minimum width=3.8cm, minimum height=1.2cm},
    arrow/.style={-Stealth, thick, color=gray!80!black}
  ]

    \node[box] (react) {\textbf{React 19 SPA}\\Client UI \& Routing};
    \node[box, right=of react] (game) {\textbf{Game Engine}\\State \& Physics};
    \node[serverbox, below=2cm of react] (express) {\textbf{Express.js API}\\JWT Auth \& Controller};
    \node[dbbox, below=2cm of game] (db) {\textbf{SQLite Database}\\\texttt{users} Table};

    \draw[arrow] (react) -- node[above, font=\tiny] {State Updates} (game);
    \draw[arrow] (react) -- node[left, font=\tiny] {HTTP / JSON} (express);
    \draw[arrow] (game) -- node[right, font=\tiny] {POST /score (JWT)} (express);
    \draw[arrow] (express) -- node[above, font=\tiny] {SQL / Transaction} (db);

    \begin{scope}[on background layer]
      \node[draw=gray!50, dashed, fill=gray!5, rounded corners=6pt, fit=(react) (game), label=above:\textbf{Client Layer (Browser / Port 5173)}] {};
      \node[draw=gray!50, dashed, fill=gray!5, rounded corners=6pt, fit=(express) (db), label=below:\textbf{Backend Layer (Node.js / Port 5000)}] {};
    \end{scope}
  \end{tikzpicture}
  \caption{Schichtenarchitektur und Datenfluss der Gesamtanwendung}
  \label{fig:arch-diag}
\end{figure}

\subsection{Sequenzdiagramm: Authentifizierung \& Score-Übermittlung}
Der Prozess der geschützten Score-Übermittlung nach einer absolvierten Spielrunde ist in Abbildung~\ref{fig:seq-diag} als UML-Sequenzdiagramm dargestellt.

\begin{figure}[htbp]
  \centering
  \begin{tikzpicture}[
    font=\sffamily\small,
    lifeline/.style={draw=gray, dashed, thick},
    entity/.style={draw=black, fill=blue!10, rectangle, minimum width=2.5cm, minimum height=0.8cm, rounded corners=2pt},
    msg/.style={-Stealth, thick},
    ret/.style={-Stealth, dashed, thick}
  ]
    % Node definitions
    \node[entity] (client) {React Client};
    \node[entity, right=2.5cm of client] (server) {Express Server};
    \node[entity, right=2.5cm of server] (db) {SQLite DB};

    % Lifelines
    \draw[lifeline] (client) -- ++(0,-6.5);
    \draw[lifeline] (server) -- ++(0,-6.5);
    \draw[lifeline] (db) -- ++(0,-6.5);

    % Messages
    \draw[msg] ([yshift=-1cm]client.south) -- node[above, font=\tiny] {1. POST /login \{user, pass\}} ([yshift=-1cm]server.south);
    \draw[msg] ([yshift=-1.8cm]server.south) -- node[above, font=\tiny] {2. SELECT * FROM users} ([yshift=-1.8cm]db.south);
    \draw[ret] ([yshift=-2.4cm]db.south) -- node[above, font=\tiny] {3. User Data + Hash} ([yshift=-2.4cm]server.south);
    \draw[ret] ([yshift=-3.0cm]server.south) -- node[above, font=\tiny] {4. 200 OK + JWT Token} ([yshift=-3.0cm]client.south);

    \draw[msg] ([yshift=-4.2cm]client.south) -- node[above, font=\tiny] {5. POST /score \{score\} + Bearer Token} ([yshift=-4.2cm]server.south);
    \draw[msg] ([yshift=-5.0cm]server.south) -- node[above, font=\tiny] {6. UPDATE users SET highScore...} ([yshift=-5.0cm]db.south);
    \draw[ret] ([yshift=-5.6cm]db.south) -- node[above, font=\tiny] {7. Transaction Success} ([yshift=-5.6cm]server.south);
    \draw[ret] ([yshift=-6.2cm]server.south) -- node[above, font=\tiny] {8. 200 OK \{highScore\}} ([yshift=-6.2cm]client.south);

  \end{tikzpicture}
  \caption{UML-Sequenzdiagramm des Authentifizierungs- und Score-Speicherprozesses}
  \label{fig:seq-diag}
\end{figure}

\subsection{Datenmodell und Relationales Schema}
Das System nutzt ein minimalistisches, aber performantes Datenbankschema. Die Tabelle \texttt{users} speichert alle relevanten Informationen zu Registrierungen und Rekorden.

\begin{figure}[htbp]
  \centering
  \begin{tikzpicture}[
    font=\sffamily\small,
    node distance=0cm,
    entity/.style={draw=black, fill=gray!10, rectangle, minimum width=5cm}
  ]
    \node[entity, fill=blue!20] (header) {\textbf{TABLE: users}};
    \node[entity, below=of header] (f1) {\texttt{id: INTEGER (PK, AUTOINCREMENT)}};
    \node[entity, below=of f1] (f2) {\texttt{username: TEXT (UNIQUE, NOT NULL)}};
    \node[entity, below=of f2] (f3) {\texttt{password: TEXT (NOT NULL)}};
    \node[entity, below=of f3] (f4) {\texttt{highScore: INTEGER (DEFAULT 0)}};
  \end{tikzpicture}
  \caption{ERD-Darstellung der Nutzertabelle in SQLite}
  \label{fig:erd}
\end{figure}

% Kap. 4: Frontend- und Game-Engine-Implementierung
\section{Frontend- und Game-Engine-Implementierung}

\subsection{Komponentenarchitektur}
Das Frontend ist nach dem Single-Page-Application-Muster aufgebaut. Die Navigation erfolgt clientseitig ohne Neuladen der Seite über React Router v7.

\begin{itemize}
  \item \textbf{\texttt{App.tsx}:} Haupt-Einstiegspunkt. Definiert die Routen (\texttt{/}, \texttt{/login}, \texttt{/register}, \texttt{/game}, \texttt{/leaderboard}) und den Zustand für den Darkmode.
  \item \textbf{\texttt{Index.tsx}:} Landingpage mit Lore-Beschreibung, Visuals und Feature-Übersicht.
  \item \textbf{\texttt{Game.tsx}:} Herzstück der Anwendung. Enthält das Render-Canvas, die physikalische Berechnungslogik und die Anbindung an die Score-API.
  \item \textbf{\texttt{LoginForm.tsx}:} Wiederverwendbare Formular-Komponente für Registrierung und Anmeldung mit integrierter Client-Validierung.
  \item \textbf{\texttt{LeaderboardTable.tsx}:} Stellt die Top 10 Spieler in einer stilisierten Tabelle dar.
\end{itemize}

\subsection{Physikalische Spielberechnung und Game Loop}
Die Physik des Spiels wird in der Komponente \texttt{Game.tsx} über einen zeitgesteuerten Intervall-Tick simuliert.

\subsubsection{Mathematische Modellierung}
Die vertikale Position des Vogels $y_t \in [0, H_{\text{game}}]$ zum diskreten Zeitschritt $t$ errechnet sich nach der Differenzen-Gleichung:
\begin{equation}
  y_{t+1} = y_t + v_t
\end{equation}
Wobei die Fallgeschwindigkeit $v_t$ kontinuierlich durch die Beschleunigung $g = 0{,}6$ (Gravitation) erhöht wird:
\begin{equation}
  v_{t+1} = v_t + g
\end{equation}

Bei einer Benutzereingabe (Leertaste oder Mausklick) wird ein konstanter negativer Sprungimpuls injiziert:
\begin{equation}
  v_{\text{jump}} = -8{,}5
\end{equation}

\subsubsection{Hindernisgenerierung und Kollisionsdetektion}
Röhren-Hindernisse bewegen sich mit einer konstanten horizontalen Geschwindigkeit $v_{\text{pipe}} = 3\,\text{px/tick}$ nach links:
\begin{equation}
  x_{\text{pipe}, t+1} = x_{\text{pipe}, t} - v_{\text{pipe}}
\end{equation}

Eine Kollision wird ausgelöst, sobald der Bounding-Box-Test positiv ausfällt:
\begin{align}
  \text{Kollision}_{\text{Boden/Decke}} &\iff y_t \le 0 \quad \lor \quad y_t + h_{\text{bird}} \ge H_{\text{game}} \\
  \text{Kollision}_{\text{Röhre}} &\iff (x_{\text{bird}} + w_{\text{bird}} > x_{\text{pipe}}) \land (x_{\text{bird}} < x_{\text{pipe}} + w_{\text{pipe}}) \nonumber \\
  &\quad \land \left( y_{\text{bird}} < h_{\text{topPipe}} \quad \lor \quad y_{\text{bird}} + h_{\text{bird}} > h_{\text{topPipe}} + g_{\text{gap}} \right)
\end{align}
Wobei $w_{\text{bird}} = 30$, $h_{\text{bird}} = 30$, $w_{\text{pipe}} = 60$, $g_{\text{gap}} = 120$ und $H_{\text{game}} = 500$ Pixel betragen.

\subsection{Theme Switching (Light- / Darkmode)}
Das Theme-Switching wird durch einfache DOM-Klassenmanipulation gelöst. In der Anwendung schaltet die Funktion \texttt{enableDarkmode()} die Klasse \texttt{light-mode} auf dem \texttt{body}-Element um, was wiederum globale CSS-Variablen für Hintergründe und Schriftfarben anpasst:

\begin{lstlisting}[language=TypeScript, caption={Theme Switch Funktionalität}]
function enableDarkmode() {
  document.body.classList.toggle("light-mode");
}
\end{lstlisting}

% Kap. 5: Backend- und API-Implementierung
\section{Backend-API und Sicherheitsarchitektur}

\subsection{RESTful API Spezifikation}
Das Express-Backend stellt vier zentrale Endpunkte zur Verfügung.

\begin{table}[htbp]
  \centering
  \begin{tabularx}{\textwidth}{lXll}
    \toprule
    \textbf{Endpunkt} & \textbf{Funktion} & \textbf{Auth-Level} & \textbf{HTTP Status} \\ \midrule
    \texttt{POST /api/users/register} & Erstellt neues Benutzerkonto & Öffentlich & 201 / 400 / 500 \\
    \texttt{POST /api/users/login} & Validiert Login \& liefert JWT & Öffentlich & 200 / 400 / 500 \\
    \texttt{GET /api/game/leaderboard} & Bietet Top-10 Bestenliste & Öffentlich & 200 / 500 \\
    \texttt{POST /api/game/score} & Aktualisiert Highscore & JWT Bearer & 200 / 401 / 500 \\ \bottomrule
  \end{tabularx}
  \caption{Spezifikation der REST-API-Endpunkte}
  \label{tab:api-routes}
\end{table}

\subsection{JWT-Authentication Middleware}
Zur Absicherung geschützter Routen prüft die Middleware \texttt{authenticateToken} den eingehenden HTTP-Header auf Vorhandensein und Validität des Tokens:

\begin{lstlisting}[language=TypeScript, caption={JWT Verifikations-Middleware}]
const authenticateToken = (req: AuthenticatedRequest, res: Response, next: NextFunction) => {
  const authHeader = req.headers['authorization'];
  const token = authHeader && authHeader.split(' ')[1];

  if (!token) return res.status(401).json({ message: "Access Token missing." });

  jwt.verify(token, JWT_SECRET, (err, user) => {
    if (err) return res.status(403).json({ message: "Invalid Token." });
    req.user = user;
    next();
  });
};
\end{lstlisting}

\subsection{Sicherheitsbewertung nach OWASP Standards}
Das Sicherheitskonzept wurde gegen die gängigsten Bedrohungen abgesichert:
\begin{itemize}
  \item \textbf{SQL Injection (SQLite):} Durch die Verwendung von Prepared Statements in \texttt{better-sqlite3} (z.\,B. \texttt{db.prepare("SELECT * FROM users WHERE id = ?").get(userId)}) werden Benutzereingaben strikt vom SQL-Code getrennt. SQLite ist somit ausgeschlossen.
  \item \textbf{Broken Authentication:} Passwörter werden niemals im Klartext gespeichert. Der Einsatz von Bcrypt schützt Passwörter in der Datenbank.
  \item \textbf{Cross-Site Scripting (XSS):} React maskiert standardmäßig alle dynamisch gerenderten Variablen im TSX-Code, wodurch die Injektion schädlicher Skripte unterbunden wird.
\end{itemize}

% Kap. 6: Testdokumentation
\section{Qualitätsmanagement und Testdokumentation}

\subsection{Teststrategie}
Die Qualitätssicherung umfasst sowohl clientseitige Eingabevalidierungen als auch strukturierte End-to-End- und Schnittstellentests der API.

\subsection{Clientseitige Formular-Validierung}
In der Komponente \texttt{LoginForm.tsx} werden Eingaben vor dem Absenden eines HTTP-Requests geprüft:
\begin{itemize}
  \item \textbf{Mindestlänge:} Passwörter müssen eine Mindestlänge von 8 Zeichen besitzen.
  \item \textbf{Konsistenzprüfung:} Bei der Registrierung müssen das Passwort und die Wiederholung identisch sein.
  \item \textbf{Checkbox-Verifikation:} Die AGB-Bestätigung ist ein zwingendes Pflichtfeld.
\end{itemize}

\subsection{API-Schnittstellentests}
Die API-Endpunkte wurden systematisch getestet. Tabelle~\ref{tab:test-results-extended} dokumentiert die Testergebnisse.

\begin{table}[htbp]
  \centering 
  \begin{tabularx}{\textwidth}{XllX}
    \toprule
    \textbf{Szenario / Testfall} & \textbf{Erwarteter Status} & \textbf{Ergebnis} & \textbf{Anmerkung} \\ \midrule
    Registrierung gültiger User & 201 Created & Erfolgreich & Nutzer in DB angelegt \\
    Registrierung existierender Name & 400 Bad Request & Erfolgreich & Unique Constraint verhindert Duplikat \\
    Login mit korrektem Passwort & 200 OK + Token & Erfolgreich & JWT wird im Body zurückgegeben \\
    Login mit falschem Passwort & 400 Bad Request & Erfolgreich & Zugriff verwehrt \\
    Score senden ohne Auth-Header & 401 Unauthorized & Erfolgreich & Middleware blockiert Anfrage \\
    Score senden mit neuem Rekord & 200 OK + DB Update & Erfolgreich & Highscore wird überschrieben \\
    Score senden mit kleinerem Wert & 200 OK (No Change) & Erfolgreich & Highscore bleibt unverändert \\ \bottomrule
  \end{tabularx}
  \caption{Testmatrix der API-Integrationstests}
  \label{tab:test-results-extended}
\end{table}

% Kap. 7: Deployment & Betrieb
\section{Deployment, Containerisierung und Betrieb}

\subsection{Multi-Container-Orchestrierung via Docker}
Um eine identische Laufzeitumgebung für Entwicklung, Testing und Produktion zu garantieren, wurde die Anwendung containerisiert. Das Zusammenspiel von Frontend und Backend wird über \texttt{compose.yaml} gesteuert.

\begin{lstlisting}[language=yaml, caption={Docker-Compose Konfiguration (\texttt{compose.yaml})}]
services:
  backend:
    build: .
    container_name: flappy_backend
    restart: always
    ports:
      - "5000:5000"
    environment:
      - PORT=5000
      - JWT_SECRET=secret-key
    volumes:
      - sqlite_data:/app/data

volumes:
  sqlite_data:
\end{lstlisting}

\subsection{Schritt-für-Schritt Inbetriebnahmeanleitung}
Für die Inbetriebnahme der Gesamtapplikation auf einem Zielsystem müssen lediglich folgende Befehle ausgeführt werden:

\begin{lstlisting}[language=bash, caption={Inbetriebnahme-Befehle}]
# 1. Repository klonen
git clone https://github.com/Guakocius/htwg.git
cd htwg/sem5/webapp

# 2. Backend containerisiert starten und im Hintergrund ausfuehren
docker-compose up --build -d

# 3. Dependencies fuer Frontend installieren
npm install --legacy-peer-deps

# 4. Frontend starten
npm run dev
\end{lstlisting}

Nach dem erfolgreichen Build-Prozess stehen die Dienste wie folgt bereit:
\begin{itemize}
  \item \textbf{Frontend Web-Applikation:} \url{http://localhost:5173}
  \item \textbf{Backend REST-API:} \url{http://localhost:5000}
\end{itemize}

% Kap. 8: Meilenstein-Dokumentation
\section{Projektfortschritt und Meilenstein-Historie}

Die Entwicklung erstreckte sich über vier definierte Meilensteine:

\subsection{Meilenstein 1: Konzeption, Wireframing \& Repository Setup}
In der ersten Phase wurden die Anforderungen analysiert, Wireframes gezeichnet sowie das Mono-Repository aufgebaut. Technologische Weichenstellungen (React, Node, SQLite) wurden in dieser Phase verbindlich festgelegt.

\subsection{Meilenstein 2: Frontend-Entwicklung \& Game-Engine}
Der Schwerpunkt des zweiten Meilensteins lag auf der Frontend-Umsetzung. Die UI-Komponenten (Landingpage, Login, Game Canvas, Leaderboard) sowie der physikalische Game-Loop in \texttt{Game.tsx} wurden vollumfänglich entwickelt.

\subsection{Meilenstein 3: Backend-REST-API \& JWT-Auth}
Im dritten Meilenstein wurde das Express-Backend realisiert. Die Anbindung der SQLite-Datenbank, das Hashing von Passwörtern mit Bcrypt sowie die Absicherung der Highscore-Übermittlung mittels JWT standen hier im Fokus.

\subsection{Meilenstein 4: Containerisierung, QA \& Dokumentation}
Der finale Meilenstein umfasste die Containerisierung mittels Docker-Compose, die Durchführung aller Systemtests, die Erstellung der Videodemo sowie die Ausarbeitung dieser Dokumentation.

% Kap. 9: Reflexion & Ausblick
\section{Reflexion und Ausblick}

\subsection{Lessons Learned \& Erfolge}
Die durchgängige Verwendung von TypeScript im gesamten Stack hat sich als wesentlicher Erfolgsfaktor erwiesen. Typsicherheitsfehler wurden frühzeitig erkannt. Die Wahl von SQLite als eingebettete Datenbank erwies sich für das Projektvolumen als ideal, da kein komplexer Datenbank-Server administriert werden musste.

\subsection{Herausforderungen}
Eine Herausforderung war die Abstimmung des React-Komponenten-Lebenszyklus mit dem asynchronen Game-Loop. Da React-Zustandsupdates (\texttt{useState}) asynchron verarbeitet werden, mussten mutable Referenzen und funktionale State-Updates verwendet werden, um präzise Kollisionsabfragen im 24-ms-Takt zu garantieren.

\subsection{Zukünftige Erweiterungen}
Mögliche zukünftige Ausbaustufen der Anwendung sind:
\begin{enumerate}
  \item \textbf{Echtzeit-Multiplayer via WebSockets:} Integration von Socket.IO, um die Positionen anderer Spieler live auf dem eigenen Spielfeld einzublenden.
  \item \textbf{Server-Side Physics Processing:} Verlegung der vollständigen Physikberechnung auf den Server, um Client-seitige Schummeleien endgültig auszuschließen.
  \item \textbf{In-Game Item \& Skin-Shop:} Freischaltung von alternativen Charakter-Skins auf Basis der erreichten Highscores.
\end{enumerate}

% Anhang
\newpage
\appendix
\section*{Anhang: Vollständiger Quellcode}
\addcontentsline{toc}{section}{Anhang: Vollständiger Quellcode}

\section{Typ-Definitionen (\texttt{types.ts})}
\begin{lstlisting}[language=TypeScript, caption={types.ts}]
export type FormProps = {
  form: string;
};

export type User = {
  id: number;
  name: string;
  password: string;
};
\end{lstlisting}

\section{Landingpage Komponente (\texttt{Index.tsx})}
\begin{lstlisting}[language=TypeScript, caption={Index.tsx}]
import flappyImg from "../../assets/flappybird.png";

/**
 * Default index page re-routing to the landing page.
 */
export default function Index() {
  return (
    <>
      <header>
        <h1>
          Welcome to Flappy Bird --
          <span className="special-text">Reimagined!</span>
        </h1>
      </header>
      <main id="center">
        <div className="flappy">
          <img
            src={flappyImg}
            id="flappy-img"
            width="auto"
            height="auto"
            alt="a yellow pixelated bird"
          />
        </div>
      </main>

      <article>
        <h2
          style={{
            display: "flex",
            alignItems: "center",
            justifyContent: "center",
          }}
        >
          <u>Lore</u>
        </h2>

        <div>
          <p>
            Long ago, a small yellow bird lived in a quiet world above endless
            green pipes. Unlike the other birds, Flappy was never meant to stay
            in a cage or rest on the ground. It dreamed of the open sky and of
            proving that even the smallest wings could overcome impossible
            obstacles. One day, the bird escaped captivity and began its journey
            through the pipe-filled land, where every flight became a test of
            courage, timing, and endurance. With each narrow gap it passed,
            Flappy came closer to freedom and to becoming a legend among birds.
          </p>
        </div>
      </article>

      <section>
        <h2>Features</h2>
        <p>
          Flappy Bird Reimagined keeps the simple challenge of the original game
          while adding its own style, atmosphere, and expanded world.
        </p>

        <ul>
          <li>Classic arcade gameplay</li>
          <li>New visual style</li>
          <li>Lore and worldbuilding</li>
          <li>Leaderboard competition</li>
        </ul>
      </section>
      <footer>&copy; 2026</footer>
    </>
  );
}
\end{lstlisting}

\section{Game Engine Komponente (\texttt{Game.tsx})}
\begin{lstlisting}[language=TypeScript, caption={Game.tsx}]
import { useState, useEffect } from "react";
import axios from "axios";
import flappyImg from "../../assets/flappybird.png";

interface Pipe {
  x: number;
  topHeight: number;
  passed: boolean;
}

/**
 * Page for displaying the game and handling the game's logic.
 */
export default function Game() {
  const [birdPosition, setBirdPosition] = useState<number>(250);
  const [velocity, setVelocity] = useState<number>(0);
  const [score, setScore] = useState<number>(0);
  const [gameOver, setGameOver] = useState<boolean>(false);
  const [gameStarted, setGameStarted] = useState<boolean>(false);
  const [pipes, setPipes] = useState<Pipe[]>([]);

  const gameHeight = 500;
  const gameWidth = 400;
  const gravity = 0.6;
  const jumpStrength = -8.5;
  const pipeWidth = 60;
  const pipeGap = 120;
  const pipeSpeed = 3;

  const handleJump = () => {
    if (!gameStarted) {
      setGameStarted(true);
    }
    if (!gameOver) {
      setVelocity(jumpStrength);
    }
  };

  useEffect(() => {
    const handleKeyDown = (e: KeyboardEvent) => {
      if (e.code === "Space") {
        handleJump();
      }
    };
    window.addEventListener("keydown", handleKeyDown);
    return () => window.removeEventListener("keydown", handleKeyDown);
  }, [gameStarted, gameOver]);

  useEffect(() => {
    if (!gameStarted || gameOver) return;

    const gameInterval = setInterval(() => {
      setBirdPosition((prev) => {
        const nextPos = prev + velocity;
        if (nextPos >= gameHeight - 30 || nextPos <= 0) {
          endGame();
        }
        return nextPos;
      });
      setVelocity((prev) => prev + gravity);

      setPipes((prevPipes) => {
        const updatedPipes = prevPipes
          .map((pipe) => {
            const nextX = pipe.x - pipeSpeed;

            if (!pipe.passed && nextX + pipeWidth < 100) {
              setScore((s) => s + 1);
              return { ...pipe, x: nextX, passed: true };
            }
            return { ...pipe, x: nextX };
          })
          .filter((pipe) => pipe.x > -pipeWidth);

        const lastPipe = updatedPipes[updatedPipes.length - 1];
        if (!lastPipe || lastPipe.x < gameWidth - 200) {
          const randomTopHeight =
            Math.floor(Math.random() * (gameHeight - pipeGap - 100)) + 40;
          updatedPipes.push({
            x: gameWidth,
            topHeight: randomTopHeight,
            passed: false,
          });
        }

        return updatedPipes;
      });

      pipes.forEach((pipe) => {
        const birdX = 100;
        const birdY = birdPosition;
        const birdSize = 30;

        if (birdX + birdSize > pipe.x && birdX < pipe.x + pipeWidth) {
          if (
            birdY < pipe.topHeight ||
            birdY + birdSize > pipe.topHeight + pipeGap
          ) {
            endGame();
          }
        }
      });
    }, 24);

    return () => clearInterval(gameInterval);
  }, [gameStarted, gameOver, velocity, birdPosition, pipes]);

  const endGame = () => {
    setGameOver(true);
    saveScore(score);
  };

  const saveScore = async (finalScore: number) => {
    const token = sessionStorage.getItem("token");
    if (!token) return;

    try {
      await axios.post(
        "http://localhost:5000/api/game/score",
        { score: finalScore },
        { headers: { Authorization: `Bearer ${token}` } },
      );
    } catch (e) {
      console.error("Failed to save high score:", e);
    }
  };

  const restartGame = () => {
    setBirdPosition(250);
    setVelocity(0);
    setScore(0);
    setPipes([]);
    setGameOver(false);
    setGameStarted(true);
  };

  return (
    <div
      onClick={handleJump}
      style={{
        display: "flex",
        flexDirection: "column",
        alignItems: "center",
        userSelect: "none",
      }}
    >
      <header
        id="game-points"
        style={{ fontSize: "2rem", marginBottom: "10px" }}
      >
        Score: {score}
      </header>
      <main
        style={{
          position: "relative",
          width: `${gameWidth}px`,
          height: `${gameHeight}px`,
          backgroundColor: "#70c5ce",
          overflow: "hidden",
          border: "2px solid #000",
        }}
      >
        <div
          id="player-container"
          style={{
            position: "absolute",
            left: "100px",
            top: `${birdPosition}px`,
            width: "30px",
            height: "30px",
          }}
        >
          <img
            src={flappyImg}
            alt="Flappy Bird"
            style={{ width: "100%", height: "100%" }}
          />
        </div>
        <div id="pipe-container">
          {pipes.map((pipe, index) => (
            <div key={index}>
              <div
                style={{
                  position: "absolute",
                  left: `${pipe.x}px`,
                  top: 0,
                  width: `${pipeWidth}px`,
                  height: `${pipe.topHeight}px`,
                  backgroundColor: "green",
                  border: "3px solid black",
                }}
              />
              <div
                style={{
                  position: "absolute",
                  left: `${pipe.x}px`,
                  top: `${pipe.topHeight + pipeGap}px`,
                  width: `${pipeWidth}px`,
                  height: `${gameHeight - pipe.topHeight - pipeGap}px`,
                  backgroundColor: "green",
                  border: "2px solid black",
                }}
              />
            </div>
          ))}
        </div>

        {!gameStarted && (
          <div
            style={{ textAlign: "center", marginTop: "200px", color: "#fff" }}
          >
            <h2>Click or Press Space to Start</h2>
          </div>
        )}

        {gameOver && (
          <div
            style={{
              position: "absolute",
              top: "30%",
              width: "100%",
              textAlign: "center",
              backgroundColor: "rgba(0, 0, 0, 0.7)",
              color: "#fff",
              padding: "20px 0",
            }}
          >
            <h2>Game Over</h2>
            <button
              onClick={restartGame}
              style={{
                padding: "10px 20px",
                fontSize: "1rem",
                cursor: "pointer",
              }}
            >
              Play Again
            </button>
          </div>
        )}
      </main>
    </div>
  );
}
\end{lstlisting}

\section{Authentifizierungs-Seiten (\texttt{Login.tsx} \& \texttt{Register.tsx})}
\begin{lstlisting}[language=TypeScript, caption={Login.tsx}]
import LoginForm from "../components/LoginForm";

import { useNavigate } from "react-router";

/**
 * Login page for the process of logging into an already
 * existing account.
 */
export default function Login() {
  const navigate = useNavigate();

  return (
    <>
      <div>
        <button onClick={() => navigate(-1)}>go back</button>
        <LoginForm form={"Login"} />
      </div>
    </>
  );
}
\end{lstlisting}

\begin{lstlisting}[language=TypeScript, caption={Register.tsx}]
import LoginForm from "../components/LoginForm";

import { useNavigate } from "react-router";

/**
 * Register page for the application process of a new user.
 */
export default function Register() {
  const navigate = useNavigate();

  return (
    <>
      <div>
        <button onClick={() => navigate(-1)}>go back</button>
        <LoginForm form={"Register"} />
      </div>
    </>
  );
}
\end{lstlisting}

\section{Leaderboard Komponente (\texttt{Leaderboard.tsx})}
\begin{lstlisting}[language=TypeScript, caption={Leaderboard.tsx}]
import { useNavigate } from "react-router";
import LeaderboardTable from "../components/LeaderboardTable";

/**
 *
 */
export default function Leaderboard() {
  const navigate = useNavigate();

  return (
    <>
      <div>
        <button onClick={() => navigate(-1)}>go back</button>
        <LeaderboardTable />
      </div>
    </>
  );
}
\end{lstlisting}

\section{Express-Backend Server (\texttt{server.ts})}
\begin{lstlisting}[language=TypeScript, caption={Auszug aus server.ts}]
import express from "express";
import type { Request, Response, NextFunction } from "express";
import cors from "cors";
import bcrypt from "bcrypt";
import jwt from "jsonwebtoken";
import dotenv from "dotenv";
import Database from "better-sqlite3";
import fs from "fs";
import { fileURLToPath } from "url";
import path from "path";

dotenv.config();

const app = express();
const PORT = process.env.PORT || 5000;
const JWT_SECRET = process.env.JWT_SECRET || "secret-key";

app.use(
  cors({
    origin: ["http://localhost:5173", "http://127.0.0.1:5173"],
    credentials: true,
    methods: ["GET", "POST", "PUT", "DELETE", "OPTIONS"],
    allowedHeaders: ["Content-Type", "Authorization"],
  }),
);
app.use(express.json());

const __filename = fileURLToPath(import.meta.url);
const __dirname = path.dirname(__filename);

const dbDir = path.join(__dirname, "../../../data");
if (!fs.existsSync(dbDir)) {
  fs.mkdirSync(dbDir, { recursive: true });
}

const db = new Database(path.join(dbDir, "flappybird.db"));

db.exec(`
        CREATE TABLE IF NOT EXISTS users (
          id INTEGER PRIMARY KEY AUTOINCREMENT,
          username TEXT UNIQUE NOT NULL,
          password TEXT NOT NULL,
          highScore INTEGER DEFAULT 0
        )
`);

console.log("Connected to SQLite successfully.");

interface UserPayload {
  id: number;
  username: string;
}

interface AuthenticatedRequest extends Request {
  user?: UserPayload | jwt.JwtPayload;
}

const authenticateToken = (
  req: AuthenticatedRequest,
  res: Response,
  next: NextFunction,
): void => {
  const authHeader = req.headers["authorization"];
  const token = authHeader && authHeader.split(" ")[1];

  if (!token) {
    res.status(401).json({ message: "Access denied. No token provided." });
    return;
  }

  try {
    const verified = jwt.verify(token, JWT_SECRET) as UserPayload;
    req.user = verified;
    next();
  } catch (e) {
    res.status(403).json({ message: `Invalid or expired token: ${e}` });
  }
};

app.post("/api/users/register", async (req: Request, res: Response) => {
  try {
    const { username, password } = req.body;

    if (!username || !password) {
      return res
        .status(400)
        .json({ message: "Username and password are required." });
    }

    const userExists = db
      .prepare("SELECT * FROM users WHERE username = ?")
      .get(username);

    if (userExists) {
      return res.status(400).json({ message: "Username is already taken." });
    }

    const salt = await bcrypt.genSalt(10);
    const hashedPassword = await bcrypt.hash(password, salt);

    const stmt = db.prepare(
      "INSERT INTO users (username, password) VALUES (?, ?)",
    );
    stmt.run(String(username), String(hashedPassword));

    return res.status(201).json({ message: "User registration successful." });
  } catch (e) {
    console.error("Registration error:", e);
    return res.status(500).json({ message: "Server error saving user data." });
  }
});

app.post("/api/users/login", async (req: Request, res: Response) => {
  try {
    const { username, password } = req.body;

    const user = db
      .prepare("SELECT * FROM users WHERE username = ?")
      .get(username) as
      | { id: number; username: string; password: string; highScore: number }
      | undefined;

    if (!user) {
      return res.status(400).json({ message: "Invalid username or password." });
    }

    const validPassword = bcrypt.compareSync(password, user.password);
    if (!validPassword) {
      return res.status(400).json({ message: "Invalid username or password" });
    }

    const token = jwt.sign(
      { id: user.id, username: user.username },
      JWT_SECRET,
      { expiresIn: "1h" },
    );

    return res.status(200).json({
      message: "Login successful.",
      token,
      user: {
        id: user.id,
        username: user.username,
        highScore: user.highScore,
      },
    });
  } catch (e) {
    console.error("Login error:", e);
    res.status(500).json({ message: "Server error during login." });
  }
});

app.get("/api/game/leaderboard", async (_req: Request, res: Response) => {
  try {
    const leaderboard = db
      .prepare(
        "SELECT id as _id, username, highScore FROM users ORDER BY highScore DESC LIMIT 10",
      )
      .all();

    return res.status(200).json(leaderboard);
  } catch (e) {
    console.error("Leaderboard error:", e);
    res.status(500).json({ message: "Server error fetching leaderboard." });
  }
});

app.post(
  "/api/game/score",
  authenticateToken,
  async (req: AuthenticatedRequest, res: Response) => {
    try {
      const { score } = req.body;
      const userId = req.user?.id;

      if (!userId) {
        res.status(401).json({ message: "Unauthorized." });
        return;
      }

      const user = db
        .prepare("SELECT * FROM users WHERE id = ?")
        .get(userId) as
        | { id: number; username: string; highScore: number }
        | undefined;

      if (!user) {
        res.status(404).json({ message: "User not found." });
        return;
      }

      if (score > user.highScore) {
        db.prepare("UPDATE users SET highScore = ? WHERE id = ?").run(
          score,
          userId,
        );
        return res
          .status(200)
          .json({ message: "New high score!", highScore: user.highScore });
      }

      return res
        .status(200)
        .json({ message: "Score processed.", highScore: user.highScore });
    } catch (e) {
      console.error("Score error:", e);
      res.status(500).json({ message: "Server error updating score." });
    }
  },
);

app.listen(PORT, () => {
  console.log(`Backend server running on port ${PORT}`);
});
\end{lstlisting}

\section{Test-Nutzerkonten}
\begin{table}[htbp]
  \centering
  \caption{Vordefinierte Accounts für die Evaluierung}
  \begin{tabular}{lll}
    \toprule
    \textbf{Rolle} & \textbf{Benutzername} & \textbf{Passwort} \\ \midrule
    Standard-User & \texttt{testuser} & \texttt{Password123!} \\
    Leaderboard Top-Player & \texttt{player1} & \texttt{SecurePass123} \\ \bottomrule
  \end{tabular}
\end{table}
\end{document}
